% ============================================================
%  DATOS DE PORTADA (bloque de título estilo IEEE)
%
%  Ya NO es una carátula a página completa: IEEEtran compone el
%  título, los autores y la afiliación en la cabecera de la primera
%  página, y el resumen va justo debajo.
%
%  RELLENA lo que está entre llaves { }:
%    - el nombre del curso
%    - el nombre del docente
%  Y VERIFICA LA LISTA DE INTEGRANTES (ver nota más abajo).
% ============================================================

\title{\LARGE\bfseries Auditoría Socioambiental de la Mina Marlin\\
       (Montana Exploradora de Guatemala, S.\,A. / Goldcorp):\\
       Análisis del Proceso de Auditoría y Propuesta de\\
       Plan de Mejora Continua bajo el Ciclo PHVA}

% ------------------------------------------------------------
%  INTEGRANTES
%  VERIFICAR: esta lista se tomó del trabajo de Perforación y Voladura,
%  asumiendo que el grupo es el mismo. Si el grupo de este curso es
%  distinto, corrige los nombres aquí. La bitácora preveía un grupo de
%  4 o 5 integrantes; aquí hay 4.
% ------------------------------------------------------------
\author{%
  \IEEEauthorblockN{%
    \textbf{Brandon Gabriel Quispe Tapia}\IEEEauthorrefmark{1},
    \textbf{Cristian Ivan Huanca Constancia}\IEEEauthorrefmark{1},\\
    \textbf{Rudy Javier Mollocondo Mollocondo}\IEEEauthorrefmark{1} y
    \textbf{Britt Lizardo Gemio Mamani}\IEEEauthorrefmark{1}%
  }\\[4pt]
  \IEEEauthorblockA{\IEEEauthorrefmark{1}Escuela Profesional de Ingeniería de Minas,
                    Facultad de Ingeniería de Minas\\
                    Universidad Nacional del Altiplano, Puno, Perú\\
                    \textit{Curso: {NOMBRE DEL CURSO} --- Evidencia N.\textsuperscript{o} 6}\\
                    \textbf{Docente:} {Grado y nombre completo del docente}}%
  \thanks{Manuscrito elaborado en 2026. Trabajo desarrollado a partir de la
  \textit{Human Rights Assessment of Goldcorp's Marlin Mine}, evaluación
  socioambiental independiente realizada por On Common Ground Consultants Inc.
  (2010) sobre la operación de Montana Exploradora de Guatemala, S.\,A.,
  departamento de San Marcos, Guatemala.}}

\markboth{Universidad Nacional del Altiplano --- Auditoría Ambiental, 2026}%
         {Auditoría Socioambiental de la Mina Marlin}
